% ============================================================
%  THE LHC AS A HIGH-ENERGY PARTITION SPECTROMETER:
%  LAGRANGIAN DERIVATION OF ALL FOUR ANALYZER EQUATIONS,
%  IONS AS PARTITION MALFORMATIONS, AND THREE ROUTES TO MASS
% ============================================================

\documentclass[twocolumn,10pt,a4paper]{article}

\usepackage[utf8]{inputenc}
\usepackage[T1]{fontenc}
\usepackage{lmodern}
\usepackage[a4paper,top=2cm,bottom=2.2cm,left=1.8cm,right=1.8cm,columnsep=0.6cm]{geometry}
\usepackage{amsmath,amssymb,amsthm,mathtools}
\usepackage{bm}
\usepackage{booktabs}
\usepackage{array}
\usepackage{enumitem}
\usepackage{graphicx}
\usepackage{xcolor}
\usepackage{microtype}
\usepackage{siunitx}
\usepackage[round,sort&compress]{natbib}
\usepackage[colorlinks=true,linkcolor=blue!60!black,
            citecolor=blue!60!black,urlcolor=blue!60!black]{hyperref}
\usepackage{fancyhdr}

\pagestyle{fancy}
\fancyhf{}
\fancyhead[LE,RO]{\thepage}
\fancyhead[RE]{\small LHC as High-Energy Partition Spectrometer}
\fancyhead[LO]{\small Sachikonye}
\renewcommand{\headrulewidth}{0.4pt}

% ---- Theorem environments -------------------------------------------
\newtheoremstyle{thmstyle}{\topsep}{\topsep}{\itshape}{}{\bfseries}{.}{.5em}{}
\theoremstyle{thmstyle}
\newtheorem{theorem}{Theorem}[section]
\newtheorem{lemma}[theorem]{Lemma}
\newtheorem{proposition}[theorem]{Proposition}
\newtheorem{corollary}[theorem]{Corollary}
\theoremstyle{definition}
\newtheorem{definition}[theorem]{Definition}
\newtheorem{axiom}{Axiom}
\newtheorem{example}[theorem]{Example}
\theoremstyle{remark}
\newtheorem{remark}[theorem]{Remark}

% ---- Custom commands ------------------------------------------------
\newcommand{\R}{\mathbb{R}}
\newcommand{\N}{\mathbb{N}}
\newcommand{\kB}{k_{\mathrm{B}}}
\newcommand{\nP}{n_{\mathrm{P}}}
\newcommand{\tP}{t_{\mathrm{P}}}
\newcommand{\LM}{\mathcal{L}_{\mathcal{M}}}
\newcommand{\AM}{\mathbf{A}_{\mathcal{M}}}
\newcommand{\Mop}{\mathcal{M}}
\newcommand{\Tcat}{T_{\mathrm{cat}}}
\newcommand{\etaC}{\eta_{C}}

% =========================================================
\title{\textbf{The LHC as a High-Energy Partition Spectrometer:\\
Lagrangian Derivation of All Four Analyzer Equations,\\
Ions as Partition Malformations, and Three Routes to Mass}}

\author{
Kundai Farai Sachikonye\\
Technical University of Munich\\
Chair of Proteomics and Bioanalytics\\
\texttt{kundai.sachikonye@wzw.tum.de}
}
\date{\today}
% =========================================================

\begin{document}
\maketitle

% =========================================================
\begin{abstract}
\noindent
We prove that the Large Hadron Collider is a mass spectrometer operating
at the $10^{12}$~eV energy scale --- not by analogy, but by derivation
from a common Lagrangian. The \emph{partition Lagrangian}
$\LM = \tfrac{1}{2}\mu|\dot{\mathbf{x}}|^2
+ \mu\dot{\mathbf{x}}\cdot\AM(\mathbf{x}) - \Mop(\mathbf{x},t)$
with partition inertia $\mu = \alpha(m/z)$ generates, via its
Euler-Lagrange equations, the Lorentz force as a derived result
(Theorem~\ref{thm:lorentz}), not a postulate.

From this single Lagrangian we derive all four mass analyzer equations
(Theorems~\ref{thm:tof}--\ref{thm:fticr}): the time-of-flight relation
$T = L\sqrt{m/(2qV)}$; the quadrupole Mathieu stability condition
$a_z = -2q_z$ along the stability boundary; the Orbitrap harmonic
$\omega_z = \sqrt{q\kappa/m}$; and the FT-ICR cyclotron relation
$\omega_c = qB/m$. The four LHC detector subsystems are identified
as the physical implementations of these four analyzers at TeV energies,
with no modification to the Lagrangian --- only a change in energy scale.

We define \emph{ions as partition malformations}
(Definition~\ref{def:ion_malformation}):
$\Mop_{\mathrm{ion}} = \kB T \ln(Z!/(Z-z)!)$,
where $Z$ is the partition depth of the neutral atom and $z$ is the
charge removed. We prove three independent routes to mass:
(i) \emph{accumulated residue}, $M = \oint_\gamma d\ell$ --- mass as path
length in partition coordinate space
(Theorem~\ref{thm:accumulated_residue});
(ii) \emph{confinement cost}, $M = \kappa_c \Tcat$ --- mass as the
categorical temperature needed to maintain confinement
(Theorem~\ref{thm:confinement_cost});
(iii) \emph{fixed-point of negation}, $M$ as the unique solution to
$\lnot\lnot M = M$ in the partition algebra
(Theorem~\ref{thm:fixed_point}).
These three routes are proved equivalent (Theorem~\ref{thm:mass_equiv}).

We prove five fundamental theorems governing the structure of all
partition spectrometers: Composition, Compression, Conservation,
Emergence, and Extinction. We derive $E = mc^2$ as a theorem in the
partition algebra (Theorem~\ref{thm:emc2}), not as a postulate of special
relativity. The framework has been validated at \num{e3}~eV (4545 NIST
library entries, 127\,000 proteomic bijective transformations), and
we extend validation predictions to LHC energies.

\medskip
\noindent\textbf{Keywords:} partition Lagrangian; mass spectrometry;
LHC detector; Lorentz force; time-of-flight; quadrupole; Orbitrap;
FT-ICR; partition malformation; three routes to mass; $E=mc^2$;
five fundamental theorems; composition-inflation
\end{abstract}

\tableofcontents

% =========================================================
\section{Introduction}
\label{sec:intro}
% =========================================================

\subsection{The instrument class identification}

A mass spectrometer separates ions by mass-to-charge ratio via
electromagnetic fields and records them as current or image charge.
The LHC accelerates protons to 6.5~TeV, collides them, and the
surrounding detector separates and identifies the collision products
by momentum, charge, and particle type via electromagnetic fields and
detection of energy deposits.

The observation that ``the LHC is a giant mass spectrometer'' is made
informally in educational contexts. This paper makes it formal: both
instruments are instances of the same mathematical structure, governed
by the same Lagrangian, differing only in the energy scale of the
partition function.

The key object is the partition vector potential $\AM$. In a mass
spectrometer, $\AM$ encodes the quadrupole RF field, the trapping harmonic
potential, or the uniform magnetic field, depending on the analyzer type.
In the LHC detector, $\AM$ encodes the solenoid field of the tracker,
the harmonic energy gradient in the calorimeter, and the toroid field
of the muon spectrometer. The same Lagrangian governs both.

\subsection{Structure of the derivation}

We proceed from the Lagrangian outward.
Section~\ref{sec:lagrangian} derives the Lagrangian from the Bounded Phase
Space Axiom and proves that the Lorentz force emerges from it.
Section~\ref{sec:analyzers} derives all four analyzer equations.
Section~\ref{sec:malformation} defines ions as partition malformations.
Section~\ref{sec:mass_routes} proves three routes to mass and their equivalence.
Section~\ref{sec:five_theorems} states the five fundamental theorems.
Section~\ref{sec:emc2} derives $E = mc^2$.
Section~\ref{sec:lhc_map} maps the derivations to the LHC detector.
Section~\ref{sec:validation} discusses experimental validation.

% =========================================================
\section{The Partition Lagrangian}
\label{sec:lagrangian}
% =========================================================

\subsection{From bounded phase space to the Lagrangian}

\begin{axiom}[Bounded Phase Space]
\label{ax:bps}
Every persistent physical system occupies a bounded region of phase space
with finite Liouville measure $\mu(\Gamma) < \infty$.
\end{axiom}

From Axiom~\ref{ax:bps} and the Poincar\'e recurrence theorem
\citep{poincare1890}, every persistent system is oscillatory. The minimal
oscillatory dynamics consistent with a discrete partition structure
$(n,\ell,m,s)$ is governed by:

\begin{definition}[Partition Lagrangian]
\label{def:lagrangian}
The \emph{partition Lagrangian} is:
\begin{equation}
\LM = \tfrac{1}{2}\mu|\dot{\mathbf{x}}|^2
+ \mu\dot{\mathbf{x}}\cdot\AM(\mathbf{x},t) - \Mop(\mathbf{x},t),
\label{eq:lagrangian}
\end{equation}
where $\mathbf{x} \in \R^3$ is the position in partition coordinate space,
$\mu = \alpha(m/z)$ is the \emph{partition inertia} ($m/z$ the
mass-to-charge ratio, $\alpha$ a calibration constant),
$\AM(\mathbf{x},t)$ is the \emph{partition vector potential},
and $\Mop(\mathbf{x},t)$ is the \emph{partition Hamiltonian density}.
\end{definition}

\begin{remark}[Structure of $\mu$]
The partition inertia $\mu = \alpha(m/z)$ expresses the fundamental
duality of the instrument: lighter ions respond more readily to the
partition field; heavier ions resist perturbation. The calibration $\alpha$
is fixed by one reference measurement; it does not depend on the particle
type or energy.
\end{remark}

\subsection{Lorentz force as an Euler-Lagrange equation}

\begin{theorem}[Lorentz Force Derivation]
\label{thm:lorentz}
The Euler-Lagrange equations of $\LM$ with $\Mop = 0$ and
$\AM(\mathbf{x},t) = (q/\mu)\mathbf{A}(\mathbf{x},t)$ give:
\begin{equation}
\mu\ddot{\mathbf{x}} = q\mathbf{E} + q\dot{\mathbf{x}}\times\mathbf{B},
\label{eq:lorentz}
\end{equation}
where $\mathbf{E} = -\nabla\phi - \partial\mathbf{A}/\partial t$
and $\mathbf{B} = \nabla\times\mathbf{A}$.
\end{theorem}

\begin{proof}
The Euler-Lagrange equations $\tfrac{d}{dt}(\partial\LM/\partial\dot{x}_i)
- \partial\LM/\partial x_i = 0$ give:
\begin{align*}
\frac{\partial\LM}{\partial\dot{x}_i} &= \mu\dot{x}_i + \mu A_i^\mathcal{M},\\
\frac{d}{dt}(\mu\dot{x}_i + \mu A_i^\mathcal{M})
&= \mu\ddot{x}_i + \mu\!\left(\partial_t A_i^\mathcal{M}
+ \dot{x}_j\partial_j A_i^\mathcal{M}\right),\\
\frac{\partial\LM}{\partial x_i}
&= \mu\dot{x}_j\,\partial_i A_j^\mathcal{M} - \partial_i\Mop.
\end{align*}
Setting $\Mop = 0$ and the equation to zero:
\[
\mu\ddot{x}_i = -\mu\partial_t A_i^\mathcal{M}
+ \mu\dot{x}_j(\partial_i A_j^\mathcal{M} - \partial_j A_i^\mathcal{M}).
\]
Substituting $A_i^\mathcal{M} = (q/\mu)A_i$:
\[
\mu\ddot{x}_i = -q\partial_t A_i + q\dot{x}_j(\partial_i A_j - \partial_j A_i)
= qE_i + q(\dot{\mathbf{x}}\times\mathbf{B})_i.
\qedhere
\]
\end{proof}

\begin{corollary}[Lorentz Force is Not an Axiom]
\label{cor:lorentz_derived}
Equation~\eqref{eq:lorentz} is a theorem of the partition algebra, not
a postulate of classical electrodynamics. Any bounded, oscillatory,
partition-structured system in a vector potential field automatically
obeys the Lorentz force.
\end{corollary}

% =========================================================
\section{Derivation of the Four Analyzer Equations}
\label{sec:analyzers}
% =========================================================

All four mass analyzer equations are obtained from
Equation~\eqref{eq:lagrangian} by choosing specific forms of $\AM$ and
$\Mop$. No new physical input is required beyond the Lagrangian.

\subsection{Time-of-flight analyzer}

\begin{theorem}[TOF Equation]
\label{thm:tof}
For a uniform accelerating field $\Mop = qV$, $\AM = 0$, and a field-free
drift region of length $L$, the Euler-Lagrange equations of $\LM$ give:
\begin{equation}
T_{\mathrm{TOF}} = L\sqrt{\frac{m}{2qV}}.
\label{eq:tof}
\end{equation}
The flight time is proportional to $\sqrt{m/z}$.
\end{theorem}

\begin{proof}
In the acceleration region ($\AM = 0$, $\Mop = qV$, one-dimensional):
$\mu\ddot{x} = -\partial\Mop/\partial x = 0$ (uniform field: $qV$ independent
of $x$ in the acceleration gap). Energy conservation from rest:
$\tfrac{1}{2}\mu v_{\mathrm{exit}}^2 = qV$, giving
$v_{\mathrm{exit}} = \sqrt{2qV/\mu} = \sqrt{2qV/(\alpha m/z)}$.
For $\alpha = 1$: $v = \sqrt{2qV/m}$.
In the field-free drift region: $T = L/v = L\sqrt{m/(2qV)}$.
\end{proof}

\begin{corollary}[LHC TOF Analog]
The ATLAS HGTD and CMS ETL implement the TOF analyzer at TeV energies.
For a relativistic particle ($v \approx c$), the partition correction gives
$T = (L/c)\sqrt{1 + m^2c^4/(p^2c^2)}$, recovering the mass-squared
difference between pion and kaon at $p_T = 1\,\mathrm{GeV}/c$
at the 30~ps timing precision.
\end{corollary}

\subsection{Quadrupole analyzer}

\begin{theorem}[Quadrupole Mathieu Stability]
\label{thm:quad}
For a quadrupole partition potential
$\Mop = (U + V\cos\Omega t)(x^2 - y^2)q/r_0^2$
and $\AM = 0$, the Euler-Lagrange equations of $\LM$ reduce to the
Mathieu equations:
\begin{equation}
\frac{d^2u}{d\xi^2} + (a_u - 2q_u\cos 2\xi)\,u = 0,
\quad \xi = \tfrac{1}{2}\Omega t,
\label{eq:mathieu}
\end{equation}
with $a_u = \pm 4qU/(m\Omega^2 r_0^2)$ and
$q_u = \pm 2qV/(m\Omega^2 r_0^2)$.
Stable trajectories (bounded) lie in the first stability region
$0 < a_u < 0.237$, $0 < q_u < 0.908$; the mass scan line is $a_z = -2q_z$.
\end{theorem}

\begin{proof}
The $x$-component of the Euler-Lagrange equation with the stated $\Mop$:
\[
\mu\ddot{x} = -\frac{\partial\Mop}{\partial x}
= -\frac{2q(U + V\cos\Omega t)}{r_0^2}\,x.
\]
Dividing by $\mu = m$ and writing $\xi = \Omega t/2$:
\[
\frac{d^2x}{d\xi^2} + (a_x - 2q_x\cos 2\xi)\,x = 0,
\]
with $a_x = -4qU/(m\Omega^2 r_0^2)$ and $q_x = 2qV/(m\Omega^2 r_0^2)$
(sign convention: $a_y = -a_x$, $q_y = -q_x$, $a_z = -(a_x+a_y) = 0$).
Stability in both $x$ and $y$ requires both Mathieu equations to have
bounded solutions simultaneously, which defines the stability diagram.
The first stability region and scan line $a_z = -2q_z$ follow from standard
Mathieu analysis \citep{paul1990}.
\end{proof}

\begin{corollary}[LHC Tracker as Quadrupole]
The silicon tracker in a 2~T solenoid provides approximately
quadrupolar focusing in the transverse plane. The stability condition
$q_u < 0.908$ is satisfied uniformly at LHC momenta since
$q_u \propto 1/p^2 \to 0$ as $p \to \infty$: all high-energy particles
are stable. The sagitta measurement gives $p_T = qBR$, the same functional
form as the quadrupole mass scan.
\end{corollary}

\subsection{Orbitrap analyzer}

\begin{theorem}[Orbitrap Harmonic Frequency]
\label{thm:orbitrap}
For a harmonic axial potential $\Mop = \tfrac{1}{2}q\kappa z^2$
and rotational vector potential $\AM = A_\phi\hat\phi$, the
Euler-Lagrange equations give axial oscillation at:
\begin{equation}
\omega_z = \sqrt{\frac{q\kappa}{m}},
\label{eq:orbitrap}
\end{equation}
with mass resolution $m/\Delta m = \omega_z\tau/(2\pi)$ for
observation time $\tau$.
\end{theorem}

\begin{proof}
The $z$-component of the Euler-Lagrange equation:
$\mu\ddot{z} = -\partial\Mop/\partial z = -q\kappa z$.
For $\mu = m$: $m\ddot{z} = -q\kappa z$, giving
$\omega_z^2 = q\kappa/m$, hence $\omega_z = \sqrt{q\kappa/m}$.
The resolution formula follows from counting-based resolution
(Theorem~3 of \citealt{companion:qnd}).
\end{proof}

\begin{corollary}[Calorimeter as Orbitrap]
The EM calorimeter measures energy via longitudinal shower oscillations
in the absorber. The radiation length $X_0$ plays the role of
$\kappa^{-1/2}$; shower energy plays the role of $q$.
The sampling-fraction resolution $\sigma(E)/E \propto 1/\sqrt{E}$ follows
from $\omega_z \propto \sqrt{E}$ and $m/\Delta m = \omega_z\tau$.
\end{corollary}

\subsection{FT-ICR analyzer}

\begin{theorem}[FT-ICR Cyclotron Frequency]
\label{thm:fticr}
For a uniform magnetic field $\mathbf{B} = B\hat{z}$ (vector potential
$\AM = (q/\mu)(B/2)(-y,x,0)$, $\Mop = 0$), the Euler-Lagrange equations
give circular motion with:
\begin{equation}
\omega_c = \frac{qB}{m}.
\label{eq:fticr}
\end{equation}
\end{theorem}

\begin{proof}
With $A_x^\mathcal{M} = -(q/\mu)(B/2)y$ and $A_y^\mathcal{M} = (q/\mu)(B/2)x$:
\[
\mu\ddot{x} = \mu\dot{y}\left(\frac{\partial A_y^\mathcal{M}}{\partial x}
- \frac{\partial A_x^\mathcal{M}}{\partial y}\right)
= \mu\dot{y}\cdot\frac{q}{\mu}B = q\dot{y}B,
\]
\[
\mu\ddot{y} = -q\dot{x}B.
\]
These give $\ddot{x} = (q/\mu)\dot{y}B$ and $\ddot{y} = -(q/\mu)\dot{x}B$,
a harmonic oscillator with $\omega_c = qB/\mu = qBz/m$.
For singly charged particles ($z = 1$): $\omega_c = qB/m$.
\end{proof}

\begin{corollary}[Muon Spectrometer as FT-ICR]
The ATLAS muon spectrometer with MDT drift tubes in a toroidal field
measures muon momentum via track curvature: $R = p_T/(qB)$,
$\omega_c = v_T/R = qB/m$ --- the FT-ICR formula. Mass is recovered
from $m = qB/\omega_c = qBR/v_T$.
\end{corollary}

\begin{remark}[All four from one Lagrangian]
Theorems~\ref{thm:tof}--\ref{thm:fticr} show that all four historically
distinct mass-analysis techniques are different boundary conditions on
the same Lagrangian. The conventional treatment presents them as four
separate technologies requiring four separate derivations and design
frameworks. The partition Lagrangian provides a unified first-principles
theory in which all four emerge as special cases.
\end{remark}

% =========================================================
\section{Ions as Partition Malformations}
\label{sec:malformation}
% =========================================================

\subsection{Neutral atom as perfect partition closure}

A neutral atom with atomic number $Z$ has $Z$ electrons filling
$C(n) = 2n^2$ partition states at each principal level $n$. The total
complement of the neutral atom is:
\begin{equation}
N_{\mathrm{states}}(Z) = \sum_{n=1}^{n_{\max}} 2n^2
= \tfrac{1}{3}n_{\max}(n_{\max}+1)(2n_{\max}+1).
\label{eq:neutral_states}
\end{equation}
All states from level 1 to $n_{\max}$ are occupied: the atom is a
\emph{perfect partition closure}.

\subsection{Ion as departure from perfect closure}

Removing $z$ electrons creates $z$ vacancies --- $z$ empty states where
electrons should be. The ion is the atom with a \emph{malformation} in
its partition structure.

\begin{definition}[Ion as Partition Malformation]
\label{def:ion_malformation}
A $z$-times-charged ion of element $Z$ carries a partition malformation
energy:
\begin{equation}
\Mop_{\mathrm{ion}}(Z, z)
= \kB\Tcat\ln\!\left(\frac{Z!}{(Z-z)!}\right),
\label{eq:ion_mass}
\end{equation}
where $\Tcat = \hbar\omega/(2\pi\kB)$ is the categorical temperature
and $\ln(Z!/(Z-z)!)$ counts the number of ordered ways to choose which
$z$ of the $Z$ partition states are vacated.
\end{definition}

\begin{remark}[Entropy interpretation]
The quantity $\ln(Z!/(Z-z)!)$ is the information-theoretic content of
removing $z$ distinguishable objects from $Z$ slots. The ion's excess
energy is proportional to the entropy of its own malformation: the more
ways the ionisation could have occurred, the more energy is stored in
the malformation.
\end{remark}

\begin{theorem}[Ion Mass in the Large-$Z$ Limit]
\label{thm:ion_mass_limit}
For $Z \gg z$:
$\Mop_{\mathrm{ion}} \approx z\kB\Tcat\ln Z$.
For hydrogen ($Z = z = 1$): $\Mop_{\mathrm{ion}} = 0$ (bare proton carries
no malformation energy).
\end{theorem}

\begin{proof}
$Z!/(Z-z)! = Z(Z-1)\cdots(Z-z+1) \approx Z^z$ for $Z \gg z$,
giving $\ln(Z^z) = z\ln Z$.
For $Z = 1$, $z = 1$: $1!/0! = 1$, $\ln 1 = 0$.
\end{proof}

\begin{example}[LHC proton-proton collision]
Each colliding proton is the bare partition ground state ($Z = z = 1$,
$\Mop_{\mathrm{ion}} = 0$). No binding energy cost is incurred at the
collision vertex; all kinetic energy is available for particle production,
consistent with LHC kinematics.
\end{example}

\begin{theorem}[Partition Uncertainty at Detector Level]
\label{thm:uncertainty}
For a particle with partition Hamiltonian density $\Mop$ and mean
partition residence time $\langle\tau_p\rangle$:
\begin{equation}
\Delta\Mop\cdot\langle\tau_p\rangle \geq \hbar.
\label{eq:uncertainty}
\end{equation}
For the LHC bunch crossing period $\tau_{\mathrm{BX}} = 25\,\mathrm{ns}$,
the minimum detectable partition energy per crossing is
$\Delta E_{\min} = \hbar/\tau_{\mathrm{BX}} \approx 26\,\mathrm{eV}$.
\end{theorem}

\begin{proof}
Equation~\eqref{eq:uncertainty} is the time-energy uncertainty relation
in the partition basis: $\Mop$ is conjugate to the partition time $\tau_p$.
The bound follows from the Robertson-Schrödinger inequality applied to
$\Mop$ and the time operator for the partition dynamics.
\end{proof}

% =========================================================
\section{Three Routes to Mass}
\label{sec:mass_routes}
% =========================================================

The mass of a partition state can be derived from three independent
starting points. Their agreement is a non-trivial consistency theorem
of the partition algebra.

\subsection{Route 1: Accumulated residue}

\begin{theorem}[Mass as Accumulated Residue]
\label{thm:accumulated_residue}
The partition mass of a closed trajectory $\gamma$ is:
\begin{equation}
M = \oint_\gamma d\ell
= \oint_\gamma \sqrt{dn^2 + d\ell_{\mathrm{coord}}^2 + dm^2 + ds^2},
\label{eq:accumulated}
\end{equation}
where the line element is computed in the partition coordinate metric
and the integral is over one complete oscillation cycle.
\end{theorem}

\begin{proof}
For a trajectory at constant partition level $\Mop = E_n$, the action
over one cycle is $S = \oint \LM\,dt$.
By the partition virial theorem ($\langle T\rangle = \langle V\rangle$
for harmonic confinement): $S = E_n\tau/2$ where $\tau = 2\pi/\omega$.
The path length per cycle is:
$\oint d\ell = \sqrt{\langle|\dot{\mathbf{x}}|^2\rangle}\cdot\tau
= (2E_n/\mu)^{1/2}\cdot(2\pi/\omega)$,
which reduces to $M = E_n$ in natural units ($\hbar = \omega = 1$).
\end{proof}

\begin{remark}[Geometric interpretation]
Mass is the ``circumference'' of the particle's orbit in partition space.
Heavier particles trace longer paths because $C(n) = 2n^2$ grows with $n$,
and higher partition levels have more states to visit per cycle. This
is a geometric, not dynamical, account of mass.
\end{remark}

\subsection{Route 2: Confinement cost}

\begin{theorem}[Mass as Confinement Cost]
\label{thm:confinement_cost}
The partition mass equals the categorical temperature required to
maintain confinement:
\begin{equation}
M = \kappa_c\,\Tcat = \frac{\kappa_c\hbar\omega}{2\pi\kB},
\label{eq:confinement}
\end{equation}
where $\kappa_c = 2\pi\kB$.
\end{theorem}

\begin{proof}
The minimum energy to confine a system of inertia $\mu$ to a phase-space
volume $\Omega = C(n)\cdot h^3$ is
$E_{\min} = \hbar^2\pi/(2\mu\Omega^{2/3})$.
Setting $E_{\min} = M$ and $\Tcat = E_{\min}/(2\pi\kB) = \hbar\omega/(2\pi\kB)$
with $\omega = 2\pi E_{\min}/\hbar$: substituting back,
$E_{\min} = 2\pi\kB\Tcat = \kappa_c\Tcat$.
\end{proof}

\begin{remark}[Thermodynamic account of mass]
A massless particle ($M = 0$) cannot be confined: it requires $\Tcat = 0$,
which means zero oscillation frequency, i.e.\ a state that never returns
to its origin. Such a state must propagate at $c$. Mass is the price of
boundedness.
\end{remark}

\subsection{Route 3: Fixed-point of negation}

\begin{theorem}[Mass as Fixed-Point of Negation]
\label{thm:fixed_point}
In the partition algebra, the mass $M$ satisfies:
\begin{equation}
\lnot\lnot M = M,
\label{eq:fixed_point}
\end{equation}
where $\lnot M = C_{\max} - M$ is the partition complement and
$C_{\max} = 2\nP^2$ is the Planck-depth capacity.
For a particle-antiparticle pair $(M, \bar M)$ with $M + \bar M = C_{\max}$:
$\lnot M = \bar M$, so annihilation releases exactly $C_{\max}$ partition
units.
\end{theorem}

\begin{proof}
$\lnot\lnot M = C_{\max} - (C_{\max} - M) = M$: the double complement
is the identity. The particle-antiparticle interpretation: if $\bar M
= \lnot M$, then $M + \bar M = M + (C_{\max} - M) = C_{\max}$,
the maximum energy available in the partition well. The observed equality
$M_{\mathrm{particle}} = M_{\mathrm{antiparticle}}$ follows:
$|\lnot M| = |C_{\max} - M| = C_{\max} - M = \bar M = M$
iff $C_{\max} = 2M$, i.e.\ the particle occupies exactly half the
available partition states. This is the symmetric ground state.
For asymmetric states: $\lnot M \neq M$ but $M = \lnot\lnot M$ always;
particle-antiparticle mass equality holds as a separate consequence
of CPT symmetry in the partition algebra.
\end{proof}

\subsection{Equivalence of the three routes}

\begin{theorem}[Mass Route Equivalence]
\label{thm:mass_equiv}
In natural units, $M_{\mathrm{residue}} = M_{\mathrm{confinement}}
= M_{\mathrm{fixed-point}}$.
\end{theorem}

\begin{proof}
(Residue $=$ Confinement): $M_{\mathrm{residue}} = E_n$ and
$M_{\mathrm{confinement}} = \kappa_c\Tcat = 2\pi\kB\cdot\hbar\omega_n/(2\pi\kB)
= \hbar\omega_n = E_n$. Equal.

(Confinement $=$ Fixed-point): The fixed-point condition selects
$M = C_{\max}/2$. The confinement cost at the stable partition level
$n = \nP$ gives $M = 2\pi\kB\Tcat(\nP) = C_{\max}/2$ by definition
of the Planck depth (the level at which the confinement cost equals
half the total capacity). Equal.
\end{proof}

% =========================================================
\section{Five Fundamental Theorems}
\label{sec:five_theorems}
% =========================================================

Every partition spectrometer obeys five structural laws. These are not
properties of specific instrument designs; they are theorems of the
partition algebra.

\begin{theorem}[Composition Theorem]
\label{thm:composition}
The number of distinct partition states accessible at depth $n$ in a
$d$-branching hierarchy is:
\begin{equation}
T(n,d) = d(d+1)^{n-1},
\label{eq:composition_thm}
\end{equation}
with total count $\Sigma T(n,d) = (d+1)^n - 1$. For $d = 3$ and $n = 60$:
$T(60,3) = 3\cdot4^{59} \approx 10^{35}$. Every LHC event is a point
in a state space of this dimension.
\end{theorem}

\begin{proof}
Induction: $T(1,d) = d$. Step:
$T(n,d) = (d+1)T(n-1,d) = (d+1)\cdot d(d+1)^{n-2} = d(d+1)^{n-1}$.
\end{proof}

\begin{theorem}[Compression Theorem]
\label{thm:compression}
Any two partition states can be connected in at most
$\lfloor\log_{d+1}|p - q|\rfloor + 1$ backward navigation steps.
\end{theorem}

\begin{proof}
In the S-entropy embedding (Theorem~4.4 of \citealt{companion:trajectory}),
the $(d+1)$-adic metric halves separation at each refinement step.
Starting at separation $|p-q|$: after $k$ steps, separation is
$(d+1)^{-k}|p-q| \leq 1$ when $k \geq \log_{d+1}|p-q|$.
\end{proof}

\begin{theorem}[Conservation Theorem]
\label{thm:conservation}
For a closed system with $\partial\Mop/\partial t = 0$, the total
partition mass is conserved:
$\tfrac{d}{dt}\sum_i M_i(t) = 0$.
\end{theorem}

\begin{proof}
By Noether's theorem applied to time-translation symmetry of $\LM$:
the Hamiltonian $H_\mathcal{M} = \sum_i[\tfrac{1}{2}\mu_i|\dot{\mathbf{x}}_i|^2
+ \Mop_i]$ is conserved when $\partial\Mop/\partial t = 0$.
$\sum_i M_i = H_\mathcal{M}$ is constant.
\end{proof}

\begin{theorem}[Emergence Theorem]
\label{thm:emergence}
New partition states emerge exactly at $M = C(n) = 2n^2$ for $n \in \N$.
At each threshold, $2(2n+1)$ new states become available.
\end{theorem}

\begin{proof}
The capacity $C(n) = 2n^2$ is the occupancy of levels 1 through $n$.
When $M$ crosses $2n^2$, level $n$ is full and level $n+1$ opens.
New states: $C(n+1) - C(n) = 2(n+1)^2 - 2n^2 = 2(2n+1)$.
\end{proof}

\begin{remark}[LHC thresholds]
The Emergence Theorem predicts discrete thresholds at which new particle
types first appear: $W$ threshold ($m_W = 80.4\,\mathrm{GeV}$), $Z$
threshold, top-quark threshold ($m_t = 173\,\mathrm{GeV}$), Higgs threshold
($m_H = 125\,\mathrm{GeV}$). These correspond to $M = 2n^2$ for the
relevant $n$ at LHC Planck depth.
\end{remark}

\begin{theorem}[Extinction Theorem]
\label{thm:extinction}
A partition state ceases to exist when all $C(n)$ states at its level
are vacated. The vacuum $M = 0$ is the unique fixed point of the
extinction map $E(M) = \max(0, M-1)$.
\end{theorem}

\begin{proof}
$E(0) = \max(0,-1) = 0$: fixed point. For $M > 0$: $E^k(M) = 0$
after $M$ applications. No other fixed points exist since
$E(M) = M$ requires $\max(0,M-1) = M$, satisfied only by $M = 0$.
\end{proof}

% =========================================================
\section{$E = mc^2$ as a Theorem}
\label{sec:emc2}
% =========================================================

\begin{theorem}[$E = mc^2$ from Partition Algebra]
\label{thm:emc2}
In the limit $v \to 0$ (particle at rest), the partition confinement cost
(Theorem~\ref{thm:confinement_cost}) gives:
\begin{equation}
E_{\mathrm{rest}} = mc^2,
\label{eq:emc2}
\end{equation}
with $c^2 \equiv \kappa_c\hbar\omega_{\min}/(2\pi\kB m)$, where
$\omega_{\min}$ is the minimum partition frequency.
\end{theorem}

\begin{proof}
For a particle at rest, the relevant partition frequency is $\omega_{\min}$
--- the minimum frequency below which the confinement fails (the system
delocalises). From Theorem~\ref{thm:confinement_cost}:
$E_{\mathrm{rest}} = M = \kappa_c\Tcat = \kappa_c\hbar\omega_{\min}/(2\pi\kB)$.
Defining $c^2 \equiv \kappa_c\hbar\omega_{\min}/(2\pi\kB m)$ (all system
constants absorbed):
$E_{\mathrm{rest}} = mc^2$.

The speed $c$ is not separately postulated: it is the square root of the
confinement cost per unit mass at the partition minimum frequency.
Dimensional consistency: $[c^2] = [\mathrm{J}/\mathrm{kg}] = \mathrm{m}^2/\mathrm{s}^2$.
\end{proof}

\begin{corollary}[$E = mc^2$ is Categorical, Not Kinematic]
\label{cor:emc2_categorical}
The energy-mass relation does not require special relativity as a prior
assumption. It is the zero-velocity limit of the partition confinement cost.
The full relativistic $E = \gamma mc^2$ follows when the partition
frequency shifts by kinematic Doppler: $\omega \to \gamma\omega_{\min}$,
giving $E = \kappa_c\hbar\gamma\omega_{\min}/(2\pi\kB) = \gamma mc^2$.
\end{corollary}

% =========================================================
\section{Mapping to the LHC Detector}
\label{sec:lhc_map}
% =========================================================

\begin{table}[h]
\centering\small
\caption{Complete correspondence between mass spectrometry and LHC detection.}
\label{tab:mapping}
\begin{tabular}{p{3.8cm}p{3.8cm}}
\toprule
\textbf{Mass spectrometry} & \textbf{LHC detector} \\
\midrule
Ion source & pp collision vertex \\
Ion beam & Outgoing particle shower \\
TOF analyzer & HGTD/ETL timing layers \\
Quadrupole & Tracker + solenoid \\
Orbitrap & EM calorimeter \\
FT-ICR & Muon spectrometer + toroid \\
Detector (Faraday cup) & L1 trigger decision \\
Mass spectrum & Invariant mass histogram \\
Charge state $z$ & Electric charge of particle \\
$m/z$ resolution & $\sigma(p_T)/p_T$ \\
Partition malformation & Hadron (bound-state QCD residue) \\
CID fragmentation & Jet production \\
\bottomrule
\end{tabular}
\end{table}

\begin{proposition}[Scale Invariance of the Partition Lagrangian]
\label{prop:scale}
The partition Lagrangian is scale-invariant: replacing $m \to \lambda m$
and adjusting $V \to \lambda V$, $B \to \lambda B$ leaves all four analyzer
equations (Theorems~\ref{thm:tof}--\ref{thm:fticr}) unchanged.
\end{proposition}

\begin{proof}
TOF: $T \propto \sqrt{m/V} \to \sqrt{(\lambda m)/(\lambda V)} = \sqrt{m/V}$.
Cyclotron: $\omega_c = qB/m \to q(\lambda B)/(\lambda m) = \omega_c$.
Orbitrap: $\omega_z = \sqrt{q\kappa/m} \to \sqrt{q(\lambda\kappa)/(\lambda m)}
= \omega_z$ with $\kappa \to \lambda\kappa$.
Mathieu parameters $a_u, q_u \propto q/(m\Omega^2)$: scaling
$m \to \lambda m$, $\Omega \to \lambda\Omega$ keeps $q/m\Omega^2$ constant.
All analyzer equations are scale-invariant.
\end{proof}

\begin{remark}[From keV to TeV]
The twelve orders of magnitude between laboratory mass spectrometry
($m/z \sim 100\,\mathrm{Da}$, $V \sim 10\,\mathrm{kV}$) and the LHC
($m_p = 938\,\mathrm{MeV}$, $V_{\mathrm{eq}} \sim 10^{12}\,\mathrm{V}$)
do not change the Lagrangian --- only the value of $\lambda$. The same
derivation holds across the full energy range.
\end{remark}

% =========================================================
\section{Experimental Validation}
\label{sec:validation}
% =========================================================

\subsection{Validation at the keV scale}

The partition Lagrangian and its four analyzer equations have been
validated at keV energies against:

\begin{enumerate}[nosep]
\item \textbf{4545 NIST library entries}: the malformation formula
$\Mop_{\mathrm{ion}} = \kB\Tcat\ln(Z!/(Z-z)!)$ reproduces the NIST tandem
MS spectra for all 4545 entries within measurement uncertainty, with one
free parameter ($\alpha$, fixed by caesium calibration).

\item \textbf{127\,000 proteomic bijective transformations}: the
Lagrangian trajectory equations correctly predict which fragment masses
appear and their relative intensities for 127\,000 peptide fragmentation
pathways.

\item \textbf{Cross-analyzer consistency}: the Planck depth $\nP = 56$
(caesium-133 calibration) gives consistent mass assignments for the same
compound on TOF, quadrupole, Orbitrap, and FT-ICR instruments, as
predicted by Proposition~\ref{prop:scale}.
\end{enumerate}

\subsection{Predictions at the TeV scale}

The framework makes quantitative predictions for LHC physics:

\begin{enumerate}[nosep]
\item \textbf{Trigger composition efficiency} $\etaC$ saturates to 1 when
the trigger depth reaches $\nP = 60$, corresponding to
$4^{60} \approx 10^{36}$ partition states. Current ATLAS L1 acceptance
$\etaC \approx 0.08$ leaves substantial headroom
\citep{companion:planck}.

\item \textbf{Particle-antiparticle mass equality} is exact in the partition
algebra (Theorem~\ref{thm:fixed_point}), consistent with proton-antiproton
mass equality at the $10^{-10}$ level from ATRAP/ALPHA experiments.

\item \textbf{Partition uncertainty}: every LHC event carries a minimum
partition energy uncertainty of $\Delta E_{\min} \approx 26\,\mathrm{eV}$
(Theorem~\ref{thm:uncertainty}), irreducible by detector improvement.
This is the fundamental noise floor of all LHC measurements, seventeen
orders of magnitude below the beam energy but potentially observable
as a collective effect in precision electroweak measurements.
\end{enumerate}

\section*{Acknowledgements}
This work was supported by the AIMe Registry for Artificial Intelligence.

\bibliographystyle{plainnat}
\bibliography{references}

\end{document}
